\documentclass[../../main.tex]{subfiles}
\begin{document}

{\bf [解]}

題意の二次方程式を$f(x) = x^2+ax+b, \, g(x) = x^2+cx+d$ とおく．
共通接線を $l : y = h(x)$ とすると
\begin{align}
\begin{cases}
f(x) - h(x) = (x-p)^2 \\[5pt]
g(x) - h(x) = (x-q)^2
\end{cases}\label{eq:1}
\quad \dots \text{①}
\end{align}
を満たすから，辺々引いて
\begin{align} 
  f(x) - g(x) = (x-p)^2 - (x-q)^2 = (2x - p - q)(q - p)
\end{align}
となり, $y = f(x)$ と $y = g(x)$ の交点の $x$ 座標は $x = \dfrac{p+q}{2}$ である．

\begin{figure}[htb]
\begin{center}
\begin{tikzpicture}[scale=1.5, >=stealth]
  % -------------------------------------------------------------
  % パラメータ設定 (接線 l: y = 0.5*x + 0.2)
  % f(x) = (x-p)^2 + h(x)
  % g(x) = (x-q)^2 + h(x)
  % p = 0.5, q = 2.5, 交点 mid = 1.5
  % -------------------------------------------------------------
  \def\P{0.5}
  \def\Q{2.5}
  \def\Mid{1.5} % (p+q)/2

  % 共通接線 h(x) の定義: h(x) = 0.5*x + 0.2
  % f(x) = (x - 0.5)^2 + 0.5*x + 0.2
  % g(x) = (x - 2.5)^2 + 0.5*x + 0.2

  % -------------------------------------------------------------
  % 面積 S の斜線塗りつぶし (ハッチング)
  % -------------------------------------------------------------
  % 1. p <= x <= (p+q)/2 : f(x) と h(x) の間
  \fill[pattern=north east lines, pattern color=gray!70]
      plot[domain=\P:\Mid, samples=50] (\x, {(\x-\P)*(\x-\P) + 0.5*\x + 0.2})
      -- plot[domain=\Mid:\P, samples=2] (\x, {0.5*\x + 0.2})
      -- cycle;

  % 2. (p+q)/2 <= x <= q : g(x) と h(x) の間
  \fill[pattern=north east lines, pattern color=gray!70]
      plot[domain=\Mid:\Q, samples=50] (\x, {(\x-\Q)*(\x-\Q) + 0.5*\x + 0.2})
      -- plot[domain=\Q:\Mid, samples=2] (\x, {0.5*\x + 0.2})
      -- cycle;

  % -------------------------------------------------------------
  % 座標軸
  % -------------------------------------------------------------
  \draw[->] (-0.5, 0) -- (3.2, 0) node[right] {$x$};
  \draw[->] (0, -0.5) -- (0, 3.2) node[above] {$y$};
  \node[below left, font=\small] at (0,0) {$O$};

  % -------------------------------------------------------------
  % 共通接線 l: y = h(x) (斜めの直線)
  % -------------------------------------------------------------
  \draw[thick, green!50!black, domain=-0.2:3.0, samples=2]
      plot (\x, {0.5*\x + 0.2}) node[above left, font=\small] {$l : y = h(x)$};

  % -------------------------------------------------------------
  % 放物線 y = f(x), y = g(x)
  % -------------------------------------------------------------
  % y = f(x)
  \draw[thick, blue!80!black, domain=-0.2:2.0, samples=100]
      plot (\x, {(\x-\P)*(\x-\P) + 0.5*\x + 0.2}) node[above left, font=\small] {$y = f(x)$};

  % y = g(x)
  \draw[thick, red!80!black, domain=1.0:3.1, samples=100]
      plot (\x, {(\x-\Q)*(\x-\Q) + 0.5*\x + 0.2}) node[above right, font=\small] {$y = g(x)$};

  % -------------------------------------------------------------
  % 接点・交点・補助線・ラベル
  % -------------------------------------------------------------
  % 接点 (p, h(p))
  \coordinate (TouchP) at (\P, {0.5*\P + 0.2});
  \fill[purple] (TouchP) circle (1.5pt);
  \draw[dashed, gray!80] (\P, 0) node[below, font=\small, black] {$p$} -- (TouchP);

  % 接点 (q, h(q))
  \coordinate (TouchQ) at (\Q, {0.5*\Q + 0.2});
  \fill[purple] (TouchQ) circle (1.5pt);
  \draw[dashed, gray!80] (\Q, 0) node[below, font=\small, black] {$q$} -- (TouchQ);

  % 放物線同士の交点 ( (p+q)/2, f((p+q)/2) )
  \coordinate (Intersection) at (\Mid, {(\Mid-\P)*(\Mid-\P) + 0.5*\Mid + 0.2});
  \fill[purple] (Intersection) circle (1.5pt);

  % 交点からの垂直破線
  \draw[dashed, gray!80] (\Mid, 0) node[below, font=\small, black] {$\frac{p+q}{2}$} -- (Intersection);

  % 領域ラベル S
  \node[font=\small, purple!80!black, fill=white, inner sep=1pt] at (1.5, 1.2) {$S$};

\end{tikzpicture}
\end{center}
\caption{2つの放物線と共通接線 $l$ で囲まれた領域 $S$}
\end{figure}

したがって, 求める面積 $S$ として \cref{eq:1} から, $p < q$ の時
\begin{align}
S 
 &= \int_p^{\frac{p+q}{2}} (x-p)^2 dx + \int_{\frac{p+q}{2}}^q (x-q)^2 dx \\
 &= \dfrac{1}{12}(q-p)^3
\end{align}
となる．$q < p$ の時も同様だから，まとめて
\begin{align}
  S = \dfrac{1}{12}|q-p|^3 
\end{align}
が求める面積である．

\end{document}
